% ============================================================
% ch05.tex —— 第 5 章 时钟算术与韩信点兵：中国剩余定理
% ============================================================
\chapter{时钟算术与韩信点兵：中国剩余定理}

第 1 章我们只是"参观"了余数王国；这一章给它配上完整的操作系统——加法、减法、乘法、解方程，一样样装上。装完之后，我们要解开一道流传一千六百年的谜题，而谜底出自一位中国将军的军营。

\section{把钟表变成数的世界}

\begin{kaiwei}
现在是 $10$ 点（$12$ 小时制）。$50$ 小时后是几点？$1000$ 小时后呢？（不许看表。）
\end{kaiwei}

$10 + 50 = 60$，$60 = 5\times12 + 0$，压扁成 $0$，是 $12$ 点（$0$ 点）。你刚在不知不觉中做了一次模 $12$ 算术。$1000$ 小时后：$1000 + 10 = 1010$，$1010 = 84\times12 + 2$，余 $2$——$2$ 点。数越大越能感到"压扁"的好处：关心的从来不是总时数，是钟面上的位置。

钟面就是一个天然的余数王国：$13$ 点就是 $1$ 点，$23$ 点就是 $11$ 点，$-2$ 点（两天前的这时）就是 $10$ 点——超过 $12$ 或者倒退到负数，都自动"绕圈"。现在把这个世界正式建起来。模 $12$ 的世界里只有 $12$ 个"数"：$0, 1, 2, \dots, 11$（用 $12$ 或用 $0$ 当"零点"是同一间房，本书统一用 $0$）。定义运算：\textbf{先按普通规则算，再把结果压扁}。例如：
\[
7 + 8 = 15 \equiv 3, \qquad 7 \times 8 = 56 \equiv 8, \qquad 5 \times 9 = 45 \equiv 9 \pmod{12}.
\]
我们把这个世界记作 $Z_{12}$。同样可以造 $Z_7$（$0$ 到 $6$）、$Z_{10}$、任何 $Z_n$。用第 1 章的替换原理，减法也没问题：$3 - 7 = -4 \equiv 8 \pmod{12}$——往前拨 $3$ 格再往回拨 $7$ 格，等于往回拨 $4$ 格，即往前拨 $8$ 格。钟表匠几百年前就会这套"负时间换正时间"。

\begin{faming}
\textbf{土名}：钟表世界、余数王国 $\to$ \textbf{正式名}：模 $n$ 剩余类环，记作 $Z_n$（大学写法 $\mathbb{Z}/n\mathbb{Z}$）。第 7 章你会明白名字里"环"字的来历：这是一个加法、乘法并肩干活的世界。
\end{faming}

\section{在钟表世界里解方程}

有了运算就能解方程——但新世界的规矩和旧世界不一样，值得逐条体验。

\paragraph{体验一：解可能不止一个。}在 $Z_{12}$ 里解 $3x \equiv 0$：试 $x = 4$，$3\times4 = 12 \equiv 0$ \checkmark；试 $x = 8$，$24 \equiv 0$ \checkmark；加上 $x = 0$，共三个解：$x \in \{0, 4, 8\}$。在普通数的世界里"$3x = 0$ 则 $x = 0$"天经地义，在钟表世界居然冒出三个解——这是新世界的第一条怪规矩。病根在哪？$3 \times 4 \equiv 0$ 说明 $3$ 和 $4$ 是一对"假零制造者"（第 1 章烧脑时刻你侦察过的现象），它们的乘积绕圈绕回了 $0$。

\paragraph{体验二：除法时有时无。}在 $Z_{12}$ 里解 $5x \equiv 1$（求 $5$ 的"倒数"）：逐个试，$5\times5 = 25 = 2\times12+1 \equiv 1$ \checkmark——解是 $x = 5$，也就是说\textbf{在这个世界里 $5$ 的倒数是它自己}。但解 $2x \equiv 1$：试遍 $0$ 到 $11$，$2x$ 的值依次是 $2,4,6,8,10,0,2,4,\dots$——全是偶数，永远轮不到 $1$。\textbf{在 $Z_{12}$ 里，$2$ 没有倒数，"除以 $2$"是非法操作}。

哪些数有倒数、哪些没有？把 $12$ 的居民挨个查一遍（拿 $12$ 的小乘法表对角线扫一眼即可）：有倒数的是 $1, 5, 7, 11$（$1^{-1}=1$，$5^{-1}=5$，$7^{-1}=7$，$11^{-1}=11$——巧了，每个都是自己的倒数）；没有倒数的是 $2,3,4,6,8,9,10$。两份名单一对比，规律跳出来：\textbf{有倒数的恰好是与 $12$ 互素的数}。这不是 $12$ 的特例，第 10 章将把这条规律证明成定理，并解释它背后站着的是第 2 章的裴蜀等式。悬念先记账。

\section{将军的难题}

现在请出主菜。《孙子算经》（成书约公元四世纪，南北朝时期）卷下第三十一题，后人称"物不知数"：

\begin{quote}\itshape
今有物不知其数，三三数之剩二，五五数之剩三，七七数之剩二。问物几何？
\end{quote}

翻译：一堆东西不知道有多少个。三个三个数，剩 $2$ 个；五个五个数，剩 $3$ 个；七个七个数，剩 $2$ 个。问这堆东西有几个？传说韩信点兵就用这个办法：不清点人数，只让士兵三次变换队形（三三、五五、七七站队），看各次剩下的零头，就能算出总兵力——几万人的军队，几十秒清点完毕。

用今天的记号写：求 $x$ 满足
\[
x \equiv 2 \pmod 3, \qquad x \equiv 3 \pmod 5, \qquad x \equiv 2 \pmod 7 .
\]

\textbf{硬办法}（先试，为了体会难度）：从满足第一个条件的数里挑：$2, 5, 8, 11, 14, 17, 20, 23, \dots$（每个加 $3$）。再筛第二个条件（除以 $5$ 余 $3$）：$8$（余 $3$ \checkmark）、$23$（余 $3$ \checkmark）、$38, 53, 68, \dots$。再筛第三个（除以 $7$ 余 $2$）：$23$（$23 = 3\times7+2$ \checkmark）——中奖！答案是 $23$。小数字还行，但要是模数换成 $97, 103, 113$，一个个筛要筛到天荒地老。

\textbf{古人的妙法}。明代程大位《算法统宗》（1592 年）把解法编成了歌诀，朗朗上口：

\begin{quote}\centering
三人同行\textbf{七十}稀，五树梅花\textbf{廿一}枝，\\
七子团圆正\textbf{半月}，除百零五便得知。
\end{quote}

歌诀的操作流程：用余数乘魔法数再相加——
\[
2\times70 + 3\times21 + 2\times15 = 140 + 63 + 30 = 233,
\]
再减去 $105$ 的倍数把数压小：$233 - 105 - 105 = 23$。答案 $23$，与硬筛一致。\textbf{$105 = 3\times5\times7$ 是三个模数的乘积}；关键是 $70, 21, 15$ 这三个"魔法数"从哪来的。查它们的身份证：

\begin{center}
\begin{tabular}{cl}
\toprule
魔法数 & 身份（除以 $3$ / $5$ / $7$ 各余几） \\
\midrule
$70$ & 余 $\mathbf{1}$；被 $5$、$7$ 整除（除尽、除尽） \\
$21$ & 被 $3$ 整除；余 $\mathbf{1}$；被 $7$ 整除 \\
$15$ & 被 $3$、$5$ 整除；余 $\mathbf{1}$ \\
\bottomrule
\end{tabular}
\end{center}

看懂了吗？每个魔法数负责一个条件：\textbf{在自己的岗位上扮演"$1$"，在别人的岗位上装死（扮演"$0"$）}。它们是怎么造出来的？以 $70$ 为例：需要"被 $5$ 和 $7$ 整除"，先乘 $5\times7 = 35$；$35$ 除以 $3$ 余 $2$（$35 = 33 + 2$），不够，试 $35$ 的倍数：$70 = 35\times2$，$70 = 69 + 1$ 除以 $3$ 余 $1$ \checkmark。找到！同理 $21 = 3\times7$（$21 = 20+1$，除以 $5$ 恰余 $1$，一步到位）；$15 = 3\times 5$（$15 = 14+1$，除以 $7$ 恰余 $1$）。"乘倍数微调余数到 $1$"这一步若一次不中，就多试几个倍数——由于模数互素，中签是必然的（这正是裴蜀等式担保的）。

\textbf{验收整套构造}。令
\[
x_0 = 2\cdot\textcolor{cheng}{70} + 3\cdot\textcolor{cheng}{21} + 2\cdot\textcolor{cheng}{15} = 233 .
\]
逐条检查：除以 $3$ 时——$70$ 贡献余数 $2\times1 = 2$；$21$ 和 $15$ 都被 $3$ 整除，装死贡献 $0$。总余数 $2 + 0 + 0 = 2$ \checkmark。除以 $5$ 时——只有 $21$ 醒着，贡献 $3\times1 = 3$；$70$、$15$ 装死。总余数 $3$ \checkmark。除以 $7$ 时——只有 $15$ 醒着，贡献 $2\times1 = 2$ \checkmark。三条全过！而减去的 $105$ 及其倍数在三个世界里都装死（$105$ 被 $3$、$5$、$7$ 同时整除），所以减完不破坏任何余数。$233 - 105 - 105 = 23$ 是最小正解；一般解是 $23 + 105k$（$k$ 任何整数）：$128, 233, 338, \dots$ 都合法。

\textbf{这就是"分而治之"}：三个条件各派一个"专职特工"（魔法数）负责，特工之间互不踩脚（别人面前装死），最后把三份情报叠加。复杂问题被拆成三个独立的小问题——这种思想在数学里叫\textbf{中国剩余定理}，在计算机里叫"混合基表示"，在你身上叫"换个角度各个击破"。

\begin{faming}
\textbf{正式名}：\textbf{中国剩余定理}（Chinese Remainder Theorem，CRT），西方文献也叫\textbf{孙子定理}——因为出处确实是《孙子算经》（此"孙子"写书在孙武之后七百年，非彼"孙子"）。一般形式：若模数 $m_1, m_2, \dots, m_k$ \textbf{两两互素}，则同余方程组
\[
x \equiv a_1 \pmod{m_1},\quad \dots,\quad x \equiv a_k \pmod{m_k}
\]
\textbf{必有解}，且解在模 $m_1m_2\cdots m_k$ 意义下\textbf{唯一}。("两两互素"即任何两个模数的最大公因数都是 $1$。)
\end{faming}

\section{条件为什么会"打架"？}

定理要求模数两两互素，这不是洁癖，看反例就懂。求一个数满足："除以 $2$ 余 $1$"且"除以 $4$ 余 $2$"。第一条说 $x$ 是\textbf{奇数}；第二条说 $x = 4k+2$ 是\textbf{偶数}。鱼与熊掌——\textbf{无解}，永远无解。病根：$2$ 和 $4$ 有公因数 $2$，两个钟面不是独立转动的，条件互相踩脚。

而 $3, 5, 7$ 两两互素，钟面各转各的，任何余数组合都能兑现。打个比方：三个互素的钟面像三个齿轮咬合精密的表盘，$105$ 小时后它们才重新对齐一次——所以每过 $105$ 个数，同样的余数组合重现一遍（这就是"解在模 $105$ 下唯一"的含义）。而 $2$ 和 $4$ 的钟面 $4$ 小时就把对方拉回原位，兼容的余数组合只有一半，剩下的组合（如"余 $1$ 与余 $2$"）物理上转不出来。

\begin{lishi}
"物不知数"的解法在秦九韶《数书九章》（1247 年）里被推广成"大衍求一术"——"求一"就是"求一个乘自己后余 $1$ 的魔法数"，与本章构造完全一致。这项工作比西方同类结果（欧拉、高斯）早了五百多年。1852 年英国传教士伟烈亚力将中国剩余定理介绍到欧洲，1876 年德国学者马蒂生发现高斯 1801 年《算术研究》中的方法与秦九韶同构，世界才公认这项成就的归属。今天，中国剩余定理是计算机科学的基础件之一：大规模并行计算里，它让机器"用几个小钟面拼一个大钟面"，绕开大数运算的溢出问题——两千年前的点兵术，正在你的 CPU 里当差。
\end{lishi}

\begin{dongshou}
\begin{itemize}
  \yisi 在 $Z_{12}$ 里解 $8x \equiv 4$。（试遍或观察：$x = 2$ 时 $16\equiv4$ \checkmark；$x = 5$ 时 $40\equiv4$ \checkmark；$x = 8$，$x = 11$ 也行——共四解。想想为什么这个方程这么多解：$\gcd(8,12) = 4$ 恰整除 $4$。这条规律第 10 章收编。）
  \yier 解"除以 $7$ 余 $3$，除以 $11$ 余 $4$"。（手工照搬魔法数流水线：找被 $11$ 整除且除以 $7$ 余 $1$ 的数——$22 = 21+1$ \checkmark；找被 $7$ 整除且除以 $11$ 余 $1$ 的数——$56 = 55+1$ \checkmark。叠加 $3\times22 + 4\times56 = 66+224 = 290$；模 $77$ 压小：$290 - 231 = 59$。验算 $59 = 8\times7+3$ \checkmark，$59 = 5\times11+4$ \checkmark。）
  \yier 用歌诀法重解"物不知数"的变体："三三数之剩一，五五数之剩二，七七数之剩三"。（$1\times70 + 2\times21 + 3\times15 = 70+42+45 = 157$；$157 - 105 = 52$。验算三条余数。）
  \yier 把 $Z_{12}$ 里"有倒数"的名单（$1,5,7,11$）和"$\gcd(\cdot, 12) = 1$"的名单对一遍，再看 $Z_{10}$（倒数名单 $1,3,7,9$）、$Z_9$（$1,2,4,5,7,8$）。三组数据支持同一个猜想——第 10 章证明它。
\end{itemize}
\end{dongshou}

\begin{shaonao}
天文学家的问题：某彗星每 $365$ 天回归一次，某行星每 $687$ 天冲日一次，今年两者同时发生，多少天后再次同时发生？（同余方程组：$x \equiv 0 \pmod{365}$，$x \equiv 0 \pmod{687}$——即 $x$ 是两者的公倍数。$365 = 5\times73$，$687 = 3\times229$，无公共素因数，互素，所以最小公倍数就是乘积 $365\times687 = 250\,755$ 天 $\approx 686.5$ 年。用第 3 章"$\gcd \times \mathrm{lcm} = ab$"核算：$\gcd(365,687)$ 用辗转相除验证确为 $1$，故 $\mathrm{lcm} = 250755$ \checkmark。下次同框要等近七百年——这就是"行星连珠"类天文事件稀罕的数学原因。）
\end{shaonao}

\begin{chengshi}
中国剩余定理对"多于两个模数"的完整证明用归纳法：两个模数的解拼好后，把"$x \equiv x_0 \pmod{m_1m_2}$"当成一个新条件与 $m_3$ 再拼，依此类推——拼法与三模数实例完全一样，没有省略任何弯。另外"两两互素"不能减弱成"总体互素"（如 $6, 10, 15$ 总体无公因数，但两两不互素，仍会打架）；歌诀中"半月"即 $15$（十五的月相是半圆，古人用月相记数）。
\end{chengshi}

钟表世界已经就绪，倒数与除法的悬念也挂好了。下一章，用它造一把锁——一把此刻正保护着全球互联网的锁，而你将全程亲手操作它。
